\begin{textsample}{POS Dim 2 – persona\_gpt – Score 58.00 – t233\_gpt.txt}  \label{ex:f2_pos_016}
Fikile Ntshangase was shot dead in her home in South Africa for resisting \textbf{coal} mining.
Berta Cáceres was murdered in Honduras after leading a campaign against a hydroelectric dam.
Marielle Franco was assassinated in Brazil after defending \textbf{human} \textbf{rights} and denouncing police brutality in Rio ’s favelas.
Nora Apique was killed in the Philippines for \textbf{standing} up against destructive projects on her ancestral \textbf{lands}.
These \textbf{women} are not only names in the past ; they are symbols of a global \textbf{pattern} of \textbf{violence} against those who dare to protect \textbf{land}, water, and \textbf{communities}.
Their \textbf{struggles} were \textbf{rooted} in \textbf{justice} : defending \textbf{territories} from extractive industries, defending \textbf{people} from militarisation and corporate \textbf{power}, and defending the \textbf{right} to participate in \textbf{decisions} that shape their \textbf{lives} and environments.
Their murders are part of an escalating \textbf{crisis} where \textbf{environmental} \textbf{destruction} and gender-based \textbf{violence} reinforce each other.
Across the \textbf{world}, \textbf{women}, girls, and \textbf{gender} minorities are at the \textbf{frontlines} of \textbf{environmental} \textbf{resistance}.
They organise \textbf{communities}, expose abuses, and confront powerful \textbf{interests}.
At the same time, they are targeted because of who they are and what they represent in \textbf{societies} shaped by sexism and patriarchal \textbf{systems}.
They \textbf{face} threats not only for opposing mines, dams, or logging, but also for \textbf{challenging} deeply \textbf{rooted} \textbf{gender} norms and \textbf{power} \textbf{structures}.
Global Witness has documented how dangerous this work has become.
In 2020 alone, 227 \textbf{land} and \textbf{environmental} \textbf{defenders} were assassinated.
One-third of those killed were Indigenous \textbf{people}, despite Indigenous \textbf{communities} making up a much smaller share of the \textbf{world} ’s population.
Their deep ties to \textbf{territory} and their \textbf{leadership} in protecting \textbf{biodiversity} put them in direct \textbf{conflict} with \textbf{governments}, \textbf{corporations}, and criminal groups seeking \textbf{profit} from \textbf{land} and \textbf{resource} \textbf{exploitation}.
Impunity is almost total : in 95 % of these cases, no one is held accountable.
This sends a clear message to perpetrators that they can silence \textbf{defenders} without consequence.
While 90 % of those killed are men, \textbf{violence} against \textbf{women} \textbf{defenders} is rising and often takes different, gendered forms. \textbf{Women} \textbf{defenders} report rape, sexual assault, and harassment, as well as blackmail and extortion that target their families and reputations.
The intention is not just to stop their activism, but to shame, isolate, and punish them for stepping outside prescribed \textbf{gender} roles.
This \textbf{violence} is unfolding in the \textbf{context} of accelerating \textbf{climate} \textbf{impacts}.
Today, 85 % of humanity is already affected by the \textbf{climate} \textbf{crisis}. \textbf{Climate} disasters, \textbf{resource} scarcity, and \textbf{displacement} intensify existing \textbf{inequalities} and make \textbf{women}, girls, and \textbf{gender} minorities even more vulnerable to abuse.
When \textbf{communities} are \textbf{pushed} to the \textbf{brink}, those who raise their \textbf{voices} against pollution, \textbf{land} grabs, and unjust \textbf{policies} often become targets.
The \textbf{struggles} of Fikile, Berta, Marielle, and Nora show how \textbf{environmental} \textbf{defense} is inseparable from the fight against sexism, racism, and colonial \textbf{legacies}.
Extractive projects and \textbf{climate} \textbf{destruction} are not gender-neutral.
They are embedded in \textbf{systems} that value \textbf{profit} over \textbf{life} and concentrate \textbf{power} in the hands of a few—often in the Global North—while the costs are borne by \textbf{communities} in the Global South, by Indigenous \textbf{peoples}, and by those already \textbf{facing} discrimination.
An ecofeminist revolution calls for confronting these intertwined \textbf{systems} of oppression.
It insists that protecting the Earth \textbf{means} dismantling patriarchal \textbf{structures} that normalise \textbf{violence} against \textbf{women} and \textbf{gender} minorities.
It \textbf{demands} that \textbf{environmental} \textbf{policies} and \textbf{climate} solutions centre the \textbf{voices} and \textbf{leadership} of those most affected, rather than silencing or sacrificing them.
This also \textbf{means} confronting the responsibility of the Global North in driving ecological \textbf{breakdown} and enabling \textbf{violence}.
Wealthy countries, \textbf{corporations}, and financial \textbf{institutions} benefit from extractive projects that \textbf{fuel} both \textbf{environmental} \textbf{destruction} and attacks on \textbf{defenders}. \textbf{Accountability} requires more than statements of concern : it requires changing trade rules, investment flows, and political alliances that currently reward those who grab \textbf{land}, pollute \textbf{ecosystems}, and repress opposition.
Protecting \textbf{environmental} defenders—especially \textbf{women}, girls, and \textbf{gender} minorities—is not a side issue ; it is \textbf{central} to any just response to the \textbf{climate} and \textbf{biodiversity} \textbf{crises}.
Legal \textbf{protection}, rapid response mechanisms, and community-based \textbf{security} are crucial, but they must go hand in hand with transforming the conditions that make this \textbf{violence} possible in the first place.
The memories of Fikile Ntshangase, Berta Cáceres, Marielle Franco, and Nora Apique call on the \textbf{world} to \textbf{act}.
Honouring their \textbf{legacies} \textbf{means} refusing to accept a \textbf{future} where defending rivers, \textbf{forests}, and neighbourhoods is a death sentence.
It \textbf{means} \textbf{standing} with those who continue their work, \textbf{demanding} \textbf{justice} in the \textbf{face} of impunity, and building an ecofeminist movement strong enough to confront both \textbf{environmental} \textbf{destruction} and gender-based \textbf{violence}, together.

% matched lemmas: accountability, act, biodiversity, breakdown, brink, central, challenge, climate, coal, community, conflict, context, corporation, crisis, decision, defender, defense, demand, destruction, displacement, ecosystem, environmental, exploitation, face, forest, frontline, fuel, future, gender, government, human, impact, inequality, institution, interest, justice, land, leadership, legacy, life, mean, pattern, people, policy, power, profit, protection, push, resistance, resource, right, root, security, society, stand, structure, struggle, system, territory, violence, voice, woman, world
\end{textsample}
